\chapter{Geltungsbereich und Grundbegriffe}
\label{b10guide:ch:foundations}

Dieser beweisfreie Leitfaden ergänzt Band~10 um zentrale Definitionen,
Charakterisierungen und Folgerungen für natürliche Zahlen. Das Format
folgt den Leitfäden zu den Bänden~3, 4, 5, 28 und~37.

\section{Gegenstand und Abgrenzung}

Grundlage sind die Peano-Eigenschaften mit vollständiger Induktion.
Die Auswahl führt über Nachfolger und Rekursion zur natürlichen
Arithmetik, Teilbarkeit und Stellenwertdarstellung. Sie beansprucht
keine vollständige Aufzählung aller bekannten oder ableitbaren Aussagen.

Allgemeine Funktionssätze einschließlich Injektivität, Surjektivität
und Bijektivität, Äquivalenzrelationen, Halb- und totale Ordnungen,
Infima und Suprema sowie ganze, rationale und reelle Zahlen bleiben
in ihren eigenen Bänden. Weitere algebraische Strukturen werden
nicht eingeführt. Funktionen und Zahlenvergleiche dienen hier nur
den konkreten Aussagen über natürliche Zahlen. Die Injektivität
des Nachfolgers gehört dabei zu den Peano-Eigenschaften.

Nicht anders bezeichnete Zahlenvariablen liegen in $\mathbb N$.
Das Präfix $10H$ kennzeichnet die Aussage-Nummern dieses Leitfadens;
die Nummerierung des Fachbands bleibt erhalten.

\section{Peano-Eigenschaften}

\begin{definition}[Natürliche Zahlen mit Null und Nachfolger]
\label{b10guide:def:peano}
Ein \emph{Peano-System} besteht aus einer Menge $N$, einem Element
$0_N\in N$ und einer Funktion $S_N\colon N\to N$ mit den Eigenschaften
\begin{enumerate}
 \item $S_N(n)\ne 0_N$ für jedes $n\in N$.
 \item Für $m,n\in N$ gilt $S_N(m)=S_N(n)\Rightarrow m=n$.
 \item Für jede Teilmenge $U\subseteq N$ gilt
 \[
  \bigl(0_N\in U\ \land\ \forall n\in U\;(S_N(n)\in U)\bigr)
  \quad\Longrightarrow\quad U=N.
 \]
\end{enumerate}
Ein fest gewähltes solches System wird als
$(\mathbb N,0,S)$ geschrieben. Seine Elemente heißen
\emph{natürliche Zahlen}; insbesondere gehört $0$ zu $\mathbb N$.
\end{definition}

Die Induktion ist hier eine Aussage über \emph{alle} Teilmengen des
Trägers. Diese Voraussetzung ist auch für die Eindeutigkeitsaussage
in Satz~\ref{b10guide:thm:categoricity} maßgeblich.

\begin{definition}[Zahlzeichen und positive natürliche Zahlen]
\label{b10guide:def:numerals}
Die ersten Zahlzeichen sind
\[
 1=S(0),\qquad 2=S(1),\qquad 3=S(2),
\]
und entsprechend weiter durch Nachfolgerbildung.
Die Menge der \emph{positiven natürlichen Zahlen} ist
$\mathbb N_+=\mathbb N\setminus\{0\}$.
\end{definition}

\section{Mengentheoretische Realisierung}

\begin{definition}[Induktive Menge]
\label{b10guide:def:inductiveset}
Eine Menge $I$ heißt \emph{induktiv}, wenn
\[
 \varnothing\in I
 \quad\text{und}\quad
 \forall x\in I\;(x\cup\{x\}\in I).
\]
\end{definition}

\begin{theorem}[Existenz der kleinsten induktiven Menge]
\label{b10guide:thm:omega}
In ZF existiert eine induktive Menge. Für jede induktive Menge $I$
ist
\[
 \omega=\bigcap\{J\subseteq I\mid J\text{ ist induktiv}\}
\]
unabhängig von der Wahl von $I$. Die Menge $\omega$ ist induktiv und
in jeder induktiven Menge enthalten. Sie ist damit die eindeutig
bestimmte kleinste induktive Menge.
\end{theorem}

\begin{theorem}[Von-Neumann-Realisierung der Peano-Eigenschaften]
\label{b10guide:thm:von-neumann}
Mit
\[
 0=\varnothing,\qquad S(n)=n\cup\{n\}\quad(n\in\omega)
\]
ist $(\omega,0,S)$ ein Peano-System.
In dieser Realisierung gilt
\[
 0=\varnothing,\qquad 1=\{0\},\qquad
 2=\{0,1\},\qquad 3=\{0,1,2\}.
\]
Für alle $n\in\omega$ gilt ferner
\[
 S(n)=n\cup\{n\},\qquad n\notin n,\qquad n\in S(n).
\]
Die Mengengleichheiten beschreiben diese konkrete Realisierung;
die nachfolgenden Peano-Aussagen setzen sie nicht voraus.
\end{theorem}
